\section{结论}\label{sec:conclusion}
本文提出粗到精多变量时间序列修复模型\Model{}，将较强上下文估计、显式多来源Patch证据与观测条件化残差精恢复结合。容量匹配配对实验识别出完整证据路径的精度贡献，且收益在精恢复后保留。公共精恢复实验区分了骨干质量与第二阶段修正的作用。结合来源、读取粒度和成本分析，结果刻画了证据组织与数值精恢复如何共同改善缺失恢复。
